\documentclass[11pt, a4paper]{article}

% bfw-style.tex — gemeinsame Style-Definitionen für alle Handouts/Aufgabenhefte
% Voraussetzung: Kompilation mit xelatex (wegen fontspec)

\usepackage[ngerman]{babel}
\usepackage{geometry}
\geometry{a4paper, top=2.2cm, bottom=2.2cm, left=2.3cm, right=2.3cm}

\usepackage{fontspec}
% Fallback-Kette: Source Sans Pro -> Calibri -> Segoe UI -> Latin Modern Sans
% (Latin Modern ist immer mit MiKTeX vorhanden)
\IfFontExistsTF{Source Sans Pro}{%
  \setmainfont{Source Sans Pro}[Ligatures=TeX]%
}{\IfFontExistsTF{Calibri}{%
  \setmainfont{Calibri}[Ligatures=TeX]%
}{\IfFontExistsTF{Segoe UI}{%
  \setmainfont{Segoe UI}[Ligatures=TeX]%
}{%
  \setmainfont{Latin Modern Sans}[Ligatures=TeX]%
}}}
\IfFontExistsTF{Source Code Pro}{%
  \setmonofont{Source Code Pro}[Scale=0.92]%
}{\IfFontExistsTF{Consolas}{%
  \setmonofont{Consolas}[Scale=0.92]%
}{%
  \setmonofont{Latin Modern Mono}[Scale=0.92]%
}}

\usepackage{xcolor}
% Farbpalette: hell, freundlich, modern
\definecolor{bfwPrimary}{HTML}{0EA5A4}    % Petrol/Türkis - Primär
\definecolor{bfwPrimaryDark}{HTML}{0F766E} % dunkler Petrol für Headings
\definecolor{bfwAccent}{HTML}{F59E0B}     % Amber - Akzent
\definecolor{bfwSuccess}{HTML}{10B981}    % Grün - Erfolg / Lösung
\definecolor{bfwWarn}{HTML}{F97316}       % Orange - Achtung
\definecolor{bfwInk}{HTML}{0F172A}        % fast Schwarz
\definecolor{bfwInkSoft}{HTML}{475569}    % gedämpftes Grau
\definecolor{bfwBgSoft}{HTML}{F8FAFC}     % off-white
\definecolor{bfwBgTeal}{HTML}{ECFEFF}     % helles Türkis
\definecolor{bfwBgAmber}{HTML}{FEF3C7}    % helles Gelb
\definecolor{bfwBgRose}{HTML}{FFE4E6}     % helles Rosé für Eselsbrücken

\usepackage[colorlinks=true, linkcolor=bfwPrimaryDark, urlcolor=bfwPrimaryDark]{hyperref}
\usepackage{enumitem}
\setlist[itemize]{leftmargin=*, itemsep=2pt, topsep=2pt}
\setlist[enumerate]{leftmargin=*, itemsep=2pt, topsep=2pt}

\usepackage[most]{tcolorbox}
\usepackage{booktabs}
\usepackage{tikz}
\usetikzlibrary{decorations.pathmorphing, positioning, shapes.geometric, arrows.meta}

% Merkkästchen — wichtige Definition / Kernpunkt
\newtcolorbox{merksatz}[1][]{%
  enhanced, breakable,
  colback=bfwBgTeal,
  colframe=bfwPrimary,
  boxrule=0.6pt,
  arc=4pt,
  left=10pt, right=10pt, top=8pt, bottom=8pt,
  fonttitle=\bfseries\color{bfwPrimaryDark},
  title={\faLightbulb\ Merksatz},
  #1
}

% Eselsbrücke — Mnemoniks
\newtcolorbox{eselsbruecke}[1][]{%
  enhanced, breakable,
  colback=bfwBgRose,
  colframe=bfwAccent,
  boxrule=0.6pt,
  arc=4pt,
  left=10pt, right=10pt, top=8pt, bottom=8pt,
  fonttitle=\bfseries\color{bfwWarn},
  title={Eselsbrücke},
  #1
}

% Beispiel-Box
\newtcolorbox{beispiel}[1][]{%
  enhanced, breakable,
  colback=bfwBgSoft,
  colframe=bfwInkSoft,
  boxrule=0.4pt,
  arc=3pt,
  left=10pt, right=10pt, top=8pt, bottom=8pt,
  fonttitle=\bfseries\color{bfwInk},
  title={Beispiel},
  #1
}

% Lösungs-Box (für Aufgabenhefte)
\newtcolorbox{loesung}[1][]{%
  enhanced, breakable,
  colback=bfwBgAmber!40,
  colframe=bfwSuccess,
  boxrule=0.6pt,
  arc=3pt,
  left=10pt, right=10pt, top=8pt, bottom=8pt,
  fonttitle=\bfseries\color{bfwSuccess!70!black},
  title={Lösung},
  #1
}

% Glossar-Eintrag
\newcommand{\glossareintrag}[2]{%
  \par\noindent\textbf{\color{bfwPrimaryDark}#1} \enspace #2\par\smallskip
}

% Section-Stil — etwas farbiger als Standard
\usepackage{titlesec}
\titleformat{\section}
  {\Large\bfseries\color{bfwPrimaryDark}}
  {\thesection}{0.6em}{}
\titleformat{\subsection}
  {\large\bfseries\color{bfwPrimary}}
  {\thesubsection}{0.5em}{}
\titleformat{\subsubsection}
  {\normalsize\bfseries\color{bfwInk}}
  {\thesubsubsection}{0.4em}{}

% TikZ-Stil "skizzenhaft" — etwas weicher als Rough.js, aber stabil
\tikzset{
  roughbox/.style = {
    draw=bfwPrimaryDark, thick, fill=bfwBgTeal,
    rounded corners=4pt, inner sep=6pt
  },
  roughaccent/.style = {
    draw=bfwAccent, thick, fill=bfwBgAmber,
    rounded corners=4pt, inner sep=6pt
  },
  roughline/.style = {
    draw=bfwInkSoft, thick, -{Stealth[length=2.5mm]}
  }
}

% Bullet-Symbole (bewusst kein FontAwesome — vermeidet Font-Installations-Probleme)
\newcommand{\faLightbulb}{$\bullet$}
\newcommand{\faBookmark}{$\bullet$}

% Header/Footer
\usepackage{fancyhdr}
\pagestyle{fancy}
\fancyhf{}
\renewcommand{\headrulewidth}{0pt}
\fancyfoot[L]{\small\color{bfwInkSoft}\thepage}
\fancyfoot[R]{\small\color{bfwInkSoft} BFW M-064 Beschaffung im Büromanagement}


\title{\vspace*{-1.5cm}\Huge\bfseries\color{bfwPrimaryDark}Aufgabenheft Tag~8\\[0.4em]
  \large\color{bfwInkSoft}\textnormal{Bezahlung von Rechnungen --- Übungen im IHK-Klausurformat}}
\author{\textsf{Beschaffung im Büromanagement}\\
\textsf{\small\copyright\ Can Siebert\,\textperiodcentered\,2026}}
\date{Selbstlernmaterial \textperiodcentered{} 21.05.2026}

\begin{document}
\maketitle
\thispagestyle{empty}

\begin{merksatz}[title={\bfseries So nutzen Sie dieses Aufgabenheft}]
Bearbeiten Sie alle Aufgaben zunächst \emph{ohne} in die Lösungen zu schauen.
Beachten Sie die \emph{Operatoren} und schreiben Sie Antworten in vollständigen Sätzen.
Lösungen ab Seite~\pageref{loesungen}.
\end{merksatz}

\tableofcontents
\vspace{1em}
\hrule

\section{Aufgaben}

\subsection*{Aufgabe~1 --- Geldfunktionen und Zahlungsarten (8 Punkte)}

\noindent\textbf{\color{bfwAccent}YouTube-Vorbereitung:}\enspace \href{https://www.youtube.com/watch?v=kPbL4EUr_1w}{Die bargeldlose Bezahlung einfach erklärt} (\url{https://www.youtube.com/watch?v=kPbL4EUr_1w})

\medskip
\begin{enumerate}[label=(\alph*)]
  \item \textbf{Nennen} Sie die drei Hauptfunktionen des Geldes. \hfill (3~P.)
  \item \textbf{Unterscheiden} Sie Bargeld und Buchgeld. \hfill (2~P.)
  \item \textbf{Ordnen} Sie folgende Zahlungen den Arten \emph{bar}, \emph{halbbar} oder \emph{bargeldlos} zu: Kassenzahlung, Bareinzahlung auf ein Konto, SEPA-Überweisung. \hfill (3~P.)
\end{enumerate}

\subsection*{Aufgabe~2 --- Bargeldlose Zahlungsformen (10 Punkte)}

\noindent\textbf{\color{bfwAccent}YouTube-Vorbereitung:}\enspace \href{https://www.youtube.com/watch?v=AxxNiSXyEzE}{SEPA-Lastschrift in 7~Schritten erklärt} (\url{https://www.youtube.com/watch?v=AxxNiSXyEzE})

\medskip
\begin{enumerate}[label=(\alph*)]
  \item \textbf{Erläutern} Sie kurz Überweisung, Lastschrift und Dauerauftrag. \hfill (3~P.)
  \item \textbf{Unterscheiden} Sie SEPA-Basis-Lastschrift und SEPA-Firmen-Lastschrift. \hfill (3~P.)
  \item \textbf{Beurteilen} Sie: Welche Zahlungsform ist für eine wiederkehrende Monatsmiete am praktischsten und warum? \hfill (4~P.)
\end{enumerate}

\subsection*{Aufgabe~3 --- Zahlungskarten und Onlinebanking (10 Punkte)}

\begin{enumerate}[label=(\alph*)]
  \item \textbf{Unterscheiden} Sie Debitkarte, Kreditkarte und Prepaid-Karte. \hfill (3~P.)
  \item \textbf{Nennen} Sie drei moderne TAN-Verfahren beim Onlinebanking. \hfill (3~P.)
  \item \textbf{Erläutern} Sie zwei Sicherheitsregeln im Onlinebanking. \hfill (4~P.)
\end{enumerate}

\subsection*{Aufgabe~4 --- Zahlungsformen im E-Commerce (10 Punkte)}

\noindent\textbf{\color{bfwAccent}YouTube-Vorbereitung:}\enspace \href{https://www.youtube.com/watch?v=GqvYviGrJZY}{SEPA-Lastschriftmandat einfach erklärt} (\url{https://www.youtube.com/watch?v=GqvYviGrJZY})

\medskip
\begin{enumerate}[label=(\alph*)]
  \item \textbf{Nennen} Sie fünf gängige Zahlungsarten im E-Commerce. \hfill (5~P.)
  \item \textbf{Beurteilen} Sie aus Käufersicht: Vor- und Nachteile von \emph{Kauf auf Rechnung} vs.\ \emph{Vorkasse}. \hfill (3~P.)
  \item \textbf{Beurteilen} Sie aus Händlersicht: Welcher Zahlungsweg minimiert das Zahlungsausfallrisiko? \hfill (2~P.)
\end{enumerate}

\subsection*{Aufgabe~5 --- Skonto rechnen (16 Punkte)}

Sie erhalten eine Rechnung über 5.000,00~€ netto. Zahlungsbedingungen: ,,Zahlbar in
30~Tagen netto Kasse oder innerhalb 14~Tagen mit 2\,\% Skonto.''

\begin{enumerate}[label=(\alph*)]
  \item \textbf{Berechnen} Sie den \emph{Skontobetrag} und den \emph{Zahlbetrag} bei Skonto-Inanspruchnahme. \hfill (4~P.)
  \item \textbf{Berechnen} Sie den \emph{effektiven Jahreszins} der Skonto-Inanspruchnahme (kaufmännisches Jahr 360~Tage). \hfill (4~P.)
  \item \textbf{Beurteilen} Sie: Lohnt sich der Skontoabzug, wenn die Hausbank einen Kontokorrent-Zinssatz von 9\,\% verlangt? Begründen Sie. \hfill (4~P.)
  \item \textbf{Erläutern} Sie, von welcher Bezugsgröße der Skonto bei diesen Konditionen abgezogen wird. \hfill (4~P.)
\end{enumerate}

\subsection*{Aufgabe~6 --- Zahlungsbedingungen interpretieren (6 Punkte)}

\begin{enumerate}[label=(\alph*)]
  \item \textbf{Erläutern} Sie folgende Zahlungsbedingung: ,,2/10 --- netto 30''. \hfill (2~P.)
  \item \textbf{Berechnen} Sie den effektiven Jahreszins für ,,2/10 --- netto 30''. \hfill (2~P.)
  \item \textbf{Nennen} Sie zwei Vorteile einer kurzen Zahlungsfrist für den Lieferanten. \hfill (2~P.)
\end{enumerate}

\vspace{1em}
\noindent\textbf{Punkte gesamt: 60}

\newpage

\section{Lösungen}\label{loesungen}

\subsection*{Lösung Aufgabe~1}

\begin{loesung}
\textbf{(a)} Drei Funktionen: Tauschmittel, Wertaufbewahrung, Recheneinheit/Wertmaßstab.

\textbf{(b)} \emph{Bargeld}: Münzen und Banknoten, gesetzliches Zahlungsmittel.
\emph{Buchgeld (Giralgeld)}: Guthaben auf Bankkonten, per Überweisung/Lastschrift/Karte verfügbar.

\textbf{(c)} Zuordnung:
\begin{itemize}[itemsep=0pt]
  \item Kassenzahlung: \emph{bar}.
  \item Bareinzahlung auf ein Konto: \emph{halbbar}.
  \item SEPA-Überweisung: \emph{bargeldlos}.
\end{itemize}
\end{loesung}

\subsection*{Lösung Aufgabe~2}

\begin{loesung}
\textbf{(a)} \emph{Überweisung}: Zahler weist seine Bank an, einen Betrag an einen Empfänger
zu transferieren. \emph{Lastschrift}: Empfänger zieht den Betrag mit Mandat des Zahlers
ein. \emph{Dauerauftrag}: vom Zahler eingerichtete wiederkehrende Überweisung mit festem
Betrag.

\textbf{(b)} \emph{SEPA-Basis-Lastschrift}: typisch B2C, 8~Wochen Erstattungsrecht ohne
Begründung. \emph{SEPA-Firmen-Lastschrift}: zwischen Unternehmen, kein Erstattungsrecht
--- höhere Sicherheit für den Empfänger.

\textbf{(c)} \emph{Dauerauftrag}: gleicher Betrag, gleicher Empfänger, fester Termin
(z.\,B.\ 1.\ jeden Monats). Vom Zahler kontrollierbar (kein Mandat an Vermieter
nötig), keine Mehrkosten, einfach in jeder Banking-App einrichtbar.
\end{loesung}

\subsection*{Lösung Aufgabe~3}

\begin{loesung}
\textbf{(a)} \emph{Debitkarte} (girocard): Sofortige Belastung des Kontos --- Deckung
muss vorhanden sein. \emph{Kreditkarte}: Verfügungsrahmen, monatliche Sammelabrechnung.
\emph{Prepaid-Karte}: vorab aufgeladenes Guthaben.

\textbf{(b)} chipTAN (Generator), photoTAN (Mobile-App fotografiert Code), pushTAN
(eigene App auf dem Smartphone). Optional: Biometrie via Smartphone.

\textbf{(c)} Beispiele:
\begin{itemize}[itemsep=0pt]
  \item Niemals TAN am Telefon weitergeben --- die Bank fragt nicht danach.
  \item Vor TAN-Eingabe immer Empfänger-IBAN und Betrag im Generator/App prüfen.
  \item Nur HTTPS-verschlüsselte Verbindungen nutzen (Schloss-Symbol).
\end{itemize}
\end{loesung}

\subsection*{Lösung Aufgabe~4}

\begin{loesung}
\textbf{(a)} Vorkasse, Kauf auf Rechnung, SEPA-Lastschrift, Kreditkarte, PayPal, Klarna,
Sofortüberweisung (fünf reichen aus).

\textbf{(b)} \emph{Kauf auf Rechnung}: + Käufer prüft Ware vor Bezahlung; -- Händler
muss in Vorleistung gehen. \emph{Vorkasse}: + Händler erhält Geld sicher; -- Käufer
trägt das Risiko, dass Ware nicht ankommt oder mangelhaft ist.

\textbf{(c)} Aus Händlersicht minimiert \emph{Vorkasse} (oder Kreditkarte mit
Autorisierung) das Zahlungsausfallrisiko, da das Geld vor Versand eingegangen ist.
\end{loesung}

\subsection*{Lösung Aufgabe~5}

\begin{loesung}
\textbf{(a)}
\begin{itemize}[itemsep=0pt]
  \item Skontobetrag $= 5\,000 \times 0{,}02 = \mathbf{100{,}00~€}$
  \item Zahlbetrag $= 5\,000 - 100 = \mathbf{4\,900{,}00~€}$
\end{itemize}

\textbf{(b)} Effektivverzinsung $= \dfrac{2 \times 360}{30 - 14} = \dfrac{720}{16} = \mathbf{45{,}00\,\%}$ p.\,a.

\textbf{(c)} Ja, der Skontoabzug lohnt sich. Der Skontoertrag entspricht einer
Effektivverzinsung von 45\,\% p.\,a.\ --- selbst bei Aufnahme eines Kontokorrent-Kredits
zu 9\,\% bleibt ein Zinsvorteil von 36~Prozentpunkten. Praxisregel: Skonto fast immer
ziehen.

\textbf{(d)} Skonto wird vom \emph{Rechnungsbetrag nach Abzug etwaiger Liefererrabatte}
berechnet. Üblicherweise vom Nettowarenwert (vor USt) oder vom Bruttowarenwert ---
im konkreten Fall sollte die Klausuraufgabe das nennen. Für unsere Aufgabe wurde der
Skonto auf den vollen Rechnungsbetrag von 5\,000~€ gerechnet.
\end{loesung}

\subsection*{Lösung Aufgabe~6}

\begin{loesung}
\textbf{(a)} ,,2/10 --- netto 30'' bedeutet: \emph{2\,\% Skonto} bei Zahlung
\emph{innerhalb 10~Tagen}, sonst voller Rechnungsbetrag fällig in \emph{30~Tagen}.

\textbf{(b)} Effektivverzinsung $= \dfrac{2 \times 360}{30 - 10} = \dfrac{720}{20} = \mathbf{36\,\%}$ p.\,a.

\textbf{(c)} Vorteile kurzer Zahlungsfristen für den Lieferanten:
\begin{itemize}[itemsep=0pt]
  \item Schnellerer Liquiditätszufluss (geringerer Finanzierungsbedarf).
  \item Geringeres Forderungsausfallrisiko.
\end{itemize}
\end{loesung}

\vfill
\hrule
\vspace{0.5em}
\noindent\small\color{bfwInkSoft}\copyright\ Can Siebert\,\textperiodcentered\,2026\,\textperiodcentered{} Beschaffung im Büromanagement\,\textperiodcentered{} Tag~8 Aufgabenheft

\end{document}
